\documentclass[UTF8,11pt]{ctexart}

\usepackage[a4paper,margin=1in]{geometry}
\usepackage{amsmath,amssymb,amsthm,mathtools}
\usepackage{bm}
\usepackage{bbm}
\usepackage{enumitem}
\usepackage{hyperref}
\usepackage{tikz-cd}

\title{Coefficient-Split 视角下 CKKS Bootstrapping Core 的子环分解}
\author{}
\date{}

% theorem environments
\newtheorem{theorem}{定理}
\newtheorem{lemma}{引理}
\newtheorem{proposition}{命题}
\newtheorem{corollary}{推论}
\theoremstyle{definition}
\newtheorem{definition}{定义}
\newtheorem{remark}{注记}
\newtheorem{example}{例子}

% commands
\newcommand{\C}{\mathbb{C}}
\newcommand{\Z}{\mathbb{Z}}
\newcommand{\RN}{R_N}
\newcommand{\Rn}{R_n}
\newcommand{\Eval}{\operatorname{Eval}}
\newcommand{\CtwoS}{\operatorname{C2S}}
\newcommand{\StwoC}{\operatorname{S2C}}
\newcommand{\BTS}{\operatorname{BTS}}
\newcommand{\Split}{\operatorname{Split}}
\newcommand{\Merge}{\operatorname{Merge}}
\newcommand{\EvalMod}{\operatorname{EvalMod}}
\newcommand{\Diag}{\operatorname{diag}}
\newcommand{\Id}{\operatorname{Id}}
\newcommand{\BlockSplit}{\operatorname{BlockSplit}}
\newcommand{\Red}{\operatorname{Red}}
\newcommand{\op}{\operatorname}

\begin{document}


\section{核心交换图（逐步版）}

为使交换结构更清晰，先约定两个记号：
\[
\mathbb{C}^{N}_{\mathrm{can}}
\quad:=\quad
\text{大环 canonical slot 空间},
\qquad
\mathbb{C}^{N}_{\mathrm{coeff}}
\quad:=\quad
\text{大环 coefficient-packed slot 空间}.
\]
类似地，
\[
\mathbb{C}^{n}_{\mathrm{can}},
\qquad
\mathbb{C}^{n}_{\mathrm{coeff}}
\]
分别表示小环的 canonical slot 空间与 coefficient-packed slot 空间。

再记
\[
S=\bigoplus_{r=0}^{k-1}S_r,
\qquad
\EvalMod_N=f^{\oplus N},
\qquad
\EvalMod_n=f^{\oplus n}.
\]

\subsection*{Step 1: \(\CtwoS\) 与 \(\Split_k\) 的交换图}

\[
\begin{tikzcd}[column sep=huge,row sep=large]
\mathbb{C}^{N}_{\mathrm{can}}
    \arrow[r,"\Split_k"]
    \arrow[d,"\CtwoS_N"']
&
\displaystyle\bigoplus_{r=0}^{k-1}\mathbb{C}^{n}_{\mathrm{can}}
    \arrow[d,"\displaystyle\bigoplus_{r=0}^{k-1}\CtwoS_n"]
\\
\mathbb{C}^{N}_{\mathrm{coeff}}
    \arrow[r,"S"']
&
\displaystyle\bigoplus_{r=0}^{k-1}\mathbb{C}^{n}_{\mathrm{coeff}}
\end{tikzcd}
\]

对应关系为
\[
\boxed{
\left(\bigoplus_{r=0}^{k-1}\CtwoS_n\right)\circ \Split_k
=
S\circ \CtwoS_N.
}
\]

\subsection*{Step 2: \(\EvalMod\) 与 coefficient 分组的交换图}

\[
\begin{tikzcd}[column sep=huge,row sep=large]
\mathbb{C}^{N}_{\mathrm{coeff}}
    \arrow[r,"S"]
    \arrow[d,"f^{\oplus N}"']
&
\displaystyle\bigoplus_{r=0}^{k-1}\mathbb{C}^{n}_{\mathrm{coeff}}
    \arrow[d,"\displaystyle\bigoplus_{r=0}^{k-1}f^{\oplus n}"]
\\
\mathbb{C}^{N}_{\mathrm{coeff}}
    \arrow[r,"S"']
&
\displaystyle\bigoplus_{r=0}^{k-1}\mathbb{C}^{n}_{\mathrm{coeff}}
\end{tikzcd}
\]

对应关系为
\[
\boxed{
\left(\bigoplus_{r=0}^{k-1}\EvalMod_n\right)\circ S
=
S\circ \EvalMod_N.
}
\]

\subsection*{Step 3: \(\StwoC\) 与 \(\Merge_k\) 的交换图}

\[
\begin{tikzcd}[column sep=huge,row sep=large]
\mathbb{C}^{N}_{\mathrm{coeff}}
    \arrow[r,"S"]
    \arrow[d,"\StwoC_N"']
&
\displaystyle\bigoplus_{r=0}^{k-1}\mathbb{C}^{n}_{\mathrm{coeff}}
    \arrow[d,"\displaystyle\bigoplus_{r=0}^{k-1}\StwoC_n"]
\\
\mathbb{C}^{N}_{\mathrm{can}}
&
\displaystyle\bigoplus_{r=0}^{k-1}\mathbb{C}^{n}_{\mathrm{can}}
    \arrow[l,"\Merge_k"']
\end{tikzcd}
\]

对应关系为
\[
\boxed{
\Merge_k\circ
\left(\bigoplus_{r=0}^{k-1}\StwoC_n\right)\circ S
=
\StwoC_N.
}
\]

\subsection*{Step 4: 将三步串起来的整体交换图}

把 Step 1, Step 2, Step 3 纵向串联，可得完整 bootstrapping core 分解图：

\[
\begin{tikzcd}[column sep=huge,row sep=large]
\mathbb{C}^{N}_{\mathrm{can}}
    \arrow[r,"\Split_k"]
    \arrow[d,"\CtwoS_N"']
&
\displaystyle\bigoplus_{r=0}^{k-1}\mathbb{C}^{n}_{\mathrm{can}}
    \arrow[d,"\displaystyle\bigoplus_{r=0}^{k-1}\CtwoS_n"]
\\
\mathbb{C}^{N}_{\mathrm{coeff}}
    \arrow[r,"S"]
    \arrow[d,"f^{\oplus N}"']
&
\displaystyle\bigoplus_{r=0}^{k-1}\mathbb{C}^{n}_{\mathrm{coeff}}
    \arrow[d,"\displaystyle\bigoplus_{r=0}^{k-1}f^{\oplus n}"]
\\
\mathbb{C}^{N}_{\mathrm{coeff}}
    \arrow[r,"S"]
    \arrow[d,"\StwoC_N"']
&
\displaystyle\bigoplus_{r=0}^{k-1}\mathbb{C}^{n}_{\mathrm{coeff}}
    \arrow[d,"\displaystyle\bigoplus_{r=0}^{k-1}\StwoC_n"]
\\
\mathbb{C}^{N}_{\mathrm{can}}
&
\displaystyle\bigoplus_{r=0}^{k-1}\mathbb{C}^{n}_{\mathrm{can}}
    \arrow[l,"\Merge_k"']
\end{tikzcd}
\]

\[
\begin{tikzcd}[column sep=huge,row sep=large]
\text{大环正常 slots}
    \arrow[r,"\text{先拆成 }k\text{ 个小环密文}"]
    \arrow[d,"\text{把大环 coefficients 放到 slots 里}"']
&
\text{\(k\) 个小环正常 slots}
    \arrow[d,"\text{各自把小环 coefficients 放到 slots 里}"]
\\
\text{大环 coefficient 向量}
    \arrow[r,"\text{按下标模 }k\text{ 分组}"']
    \arrow[d,"\text{逐坐标做 EvalMod}"']
&
\text{\(k\) 组小环 coefficient 向量}
    \arrow[d,"\text{每组逐坐标做 EvalMod}"]
\\
\text{模约减后的大环 coefficient 向量}
    \arrow[r,"\text{仍按下标模 }k\text{ 分组}"']
    \arrow[d,"\text{恢复成大环正常 slots}"']
&
\text{\(k\) 组模约减后的小环 coefficient 向量}
    \arrow[d,"\text{各自恢复成小环正常 slots}"]
\\
\text{大环正常 slots}
&
\text{\(k\) 个小环正常 slots}
    \arrow[l,"\text{合并回一个大环密文}"']
\end{tikzcd}
\]

因此整体满足
\[
\boxed{
\Merge_k
\circ
\left(
\bigoplus_{r=0}^{k-1}
\bigl(\StwoC_n\circ \EvalMod_n\circ \CtwoS_n\bigr)
\right)
\circ
\Split_k
=
\StwoC_N\circ \EvalMod_N\circ \CtwoS_N.
}
\]

若记
\[
\BTS_N:=\StwoC_N\circ \EvalMod_N\circ \CtwoS_N,
\qquad
\BTS_n:=\StwoC_n\circ \EvalMod_n\circ \CtwoS_n,
\]
则上式可简写为
\[
\boxed{
\Merge_k
\circ
\left(\bigoplus_{r=0}^{k-1}\BTS_n\right)
\circ
\Split_k
=
\BTS_N.
}
\]

\subsection*{Step 5: 考虑实现中的全局归一化因子}

若真实实现中存在全局常数 \(\gamma_k\)，则更保守地写为
\[
\boxed{
\Merge_k^{\mathrm{impl}}
\circ
\left(\bigoplus_{r=0}^{k-1}\BTS_n^{\mathrm{impl}}\right)
\circ
\Split_k^{\mathrm{impl}}
\approx
\gamma_k\,\BTS_N^{\mathrm{impl}}.
}
\]

等价地，
\[
\boxed{
\gamma_k^{-1}\,
\Merge_k^{\mathrm{impl}}
\circ
\left(\bigoplus_{r=0}^{k-1}\BTS_n^{\mathrm{impl}}\right)
\circ
\Split_k^{\mathrm{impl}}
\approx
\BTS_N^{\mathrm{impl}}.
}
\]

\end{document}
